\chapter{Conclusions}
\setcounter{section}{0}

\section{General Conclusions}

This Final Degree Project (FDP) studied whether a compact sequence representation of sign movement
can retain enough discriminative information to recognise a controlled vocabulary of isolated ASL
signs while reducing dependence on raw video appearance. The hypothesis stated in
Chapter~\ref{chap:experiments}'s preamble and traced back to the main objective in Chapter~3 can
now be answered.

The answer is \textbf{yes, within the tested vocabulary scale, with reproducible stability across
two ASL benchmarks and grounded in an explicit ethical framework}. The landmark-based GRU+Attention
pipeline reaches 63.85\% Top-1 and 86.46\% Top-5 on WLASL30 under Stratified 5-Fold
Cross-Validation, and reproduces with comparable seed-to-seed stability on MS-ASL50 (70.20\% Top-1,
90.54\% Top-5) without any dataset-specific tuning. The pipeline achieves these results while
reducing each video frame from a $224 \times 224$ RGB tensor to 63 landmark coordinates, a
2{,}389$\times$ dimensionality reduction that simultaneously discards background appearance,
reduces appearance-based bias and protects signer identity by construction.

The work also identifies the boundaries of the approach. Recognition accuracy degrades steadily as
the vocabulary grows, from 58.82\% Top-1 at 30 classes to 9.76\% at 1{,}500 classes, reflecting the
limited per-class data available in WLASL rather than a structural failure of the architecture. The
qualitative validation with 15 participants further indicates that the prototype is perceived as
useful ($4.5/5$ on perceived utility) and recommendable (Net Promoter Score $+40$), although result
comprehension remains the lowest usability dimension and a clear target for interface improvement.

Taken together, the experimental, ethical and qualitative evidence supports the project as a
research prototype that demonstrates feasibility and reproducibility at the controlled-vocabulary
scale, while explicitly delimiting the conditions under which the approach generalises.

\section{Conclusions Regarding Project Objectives}

The six specific objectives defined in Chapter~3 are addressed in the document as follows.

\paragraph{SO1: Analyse state-of-the-art techniques in computer vision and temporal
    modelling.} The State of the Art chapter (Chapter~\ref{ch:sota}) reviews convolutional
                architectures, pose estimation, MediaPipe Hands, recurrent layers, attention
                mechanisms and Transformer architectures, situating each within the SLR literature.
                The objective is therefore considered achieved.

\paragraph{SO2: Define a landmark-based data representation for ASL video clips.}
The data preparation pipeline described in Chapter~5 converts raw WLASL and MS-ASL clips into
normalised sequences of 21 three-dimensional hand landmarks per frame. The representation is
consistent across both datasets and reduces sensitivity to background, lighting and signer
appearance. The objective is achieved.

\paragraph{SO3: Design and compare candidate convolutional and temporal architectures
    for word-level ASL recognition.} Five locally tuned architecture families were compared under
                                     matched conditions on WLASL30 (Section~\ref{sec:exp_phase2}).
                                     The GRU+Attention configuration emerged as the best-performing
                                     model under the controlled comparison and was selected as the
                                     default for downstream experiments. The objective is achieved.

\paragraph{SO4: Evaluate the selected recognition pipeline through reproducible
    experimental protocols.} The selected configuration was evaluated through Stratified 5-Fold
                             Cross-Validation, multi-seed runs at three random seeds and a
                             sequence-length sweep (Sections~\ref{sec:exp_seqlen}
                             and~\ref{sec:exp_wlasl30}). The mean Top-1 of 63.85\% $\pm$ 4.50\%
                             under Cross-Validation and the seed-to-seed standard deviation of
                             2.94\,pp on the fixed split provide the reproducibility evidence
                             required for the objective. The objective is achieved.

\paragraph{SO5: Analyse the generalisation and scalability of the recognition
    pipeline.} The pipeline was retrained on MS-ASL50 without architectural changes
               (Section~\ref{sec:exp_msasl}) and on WLASL subsets from 100 to 1{,}500 classes
               (Section~\ref{sec:exp_scaling}). The cross-dataset results confirm that training
               stability is a property of the pipeline rather than of a specific benchmark, while
               the scaling experiments quantify the per-class data threshold below which
               performance degrades. The objective is achieved with the explicit limitation that
               scaling experiments beyond 30 classes are single-seed.

\paragraph{SO6: Evaluate the runtime proof of concept and the ethical implications
    of the recognition approach.} The runtime prototype is described in Chapter~5 and qualitatively
                                  validated in Section~\ref{sec:exp_qualitative} with $N=15$
                                  participants. The ethical evaluation against Google's AI
                                  Principles (Chapter~7) yields an overall maturity score of
                                  7.6/10, identifying the inference pipeline as the primary area
                                  for improvement. The objective is achieved.

\section{Next Steps}

The conclusions above identify three frontiers that extend this work, ordered by the type of
contribution they would make.

\paragraph{Technical: beyond word-level recognition and on-device miniaturisation.}
The natural technical extension is to move from isolated gloss classification to continuous ASL
sentence recognition, which remains an open frontier of sign language research. A complementary
technical track is to miniaturise the selected GRU+Attention configuration so that the model can
run entirely on a mobile device with reduced parameter count, enabling fully on-device inference
without server involvement.

\paragraph{Scientific: cross-lingual transfer and broader architectural comparison.}
The landmark representation is, by construction, language-agnostic at the geometric level: hand
anatomy is shared across sign languages, even when vocabularies are not. A scientific extension is
to transfer the pipeline from ASL to Lengua de Signos Espa\~nola (LSE), British Sign Language (BSL)
and other sign languages, evaluating whether the same architecture generalises across linguistic
boundaries. A second scientific direction, identified as a limitation in
Section~\ref{sec:exp_limitations}, is to extend the architectural comparison beyond GRU+Attention.
In particular, running the Transformer baseline of Section~\ref{sec:exp_phase2} on MS-ASL50 and on
larger WLASL subsets, using hardware capable of accommodating larger models such as more powerful
GPUs, would test whether the recurrent advantage observed at the WLASL30 scale persists in data
regimes where Transformer architectures are typically competitive.

\paragraph{Applied: privacy-first edge deployment.} The ethical analysis in Chapter~7 identified the inference pipeline as the weakest privacy point of
the current system. The applied next step is therefore to move the entire pipeline on-device,
including landmark extraction and recognition, so that only the abstract 63-dimensional vector ever
leaves the phone if any data leaves it at all. This closes the privacy gap identified in the
ethical evaluation while preserving the recognition capability validated in Part~I.

These three directions are not exclusive. A complete continuation of the work would combine
continuous-sentence modelling with cross-lingual evaluation under a privacy-first deployment
configuration, building on the foundations established in this Final Degree Project.